% Source-derived chapter 06: Zeros of $p$-adic power series
% Original source: chabauty-coleman-bound-surfaces
\chapter{Zeros of $p$-adic power series}
\label{chabauty-coleman-bound-surfaces-chapter-06}
\source{chabauty-coleman-bound-surfaces}

\begin{remark}[Source section]
\label{chabauty-coleman-bound-surfaces-chapter-06-source}
\source{chabauty-coleman-bound-surfaces}
\source{chabauty-coleman-bound-surfaces}
This chapter reproduces the corresponding section of the retrieved
LaTeX source, including its definitions, statements, examples, and
proofs. It has not yet been translated into Lean.
\notready
\end{remark}

% The source section title was: Zeros of $p$-adic power series
\label{SecZeros}
\source{chabauty-coleman-bound-surfaces}

We keep the notation from Section \ref{NotationLoc}.

%%%%
%%%%
%%%%
\subsection{One variable} Let $z$ be a variable. Given a power series $h(z)=\sum_{n=0}^\infty a_j z^j\in K[[z]]$ the radius of convergence of $h$ is defined by $\rho_h=\sup\{ r : \lim_j |a_j|r^j=0 \}$.

If $\rho_h>0$ then $h(z)$ converges on $B_1(\rho_h)$, where it defines an analytic function. In particular, the vanishing order $\ord_{z_0}(h)$ is defined at each $z_0\in B_1(\rho_h)$. For each $0<r<\rho_h$ we define $|h|_r=\max_j |a_j|r^j$, $\nu(h,r)=\max\{m : |a_m|r^m=|h|_r\}$, and 
$$
\nn_0(h,r)=\sum_{|z_0|\le r} \ord_{z_0}(h).
$$
The following classical bound will suffice for our purposes. See, for instance,  Theorem 1.21 in \cite{HuYang}.
%%
%%
\begin{lemma}
\source{chabauty-coleman-bound-surfaces}
\notready\label{LemmaZeros1var} If $h\ne 0$ and $\rho_h>0$, then for each $0<r<\rho_h$ we have $\nn_0(h,r)\le \nu(h,r)$.
\source{chabauty-coleman-bound-surfaces}
\end{lemma}
%%
%%

%%%%
%%%%
%%%%
\subsection{Several variables} Let $n\ge 1$ and let $\xx=(x_1,...,x_n)$ be an $n$-tuple of variables. Given 
$$
H(\xx)=\sum_{\alpha\in \N^n} c_\alpha \xx^\alpha\in K[[\xx ]]
$$ 
the radius of convergence $\rho_H$ is defined by $\rho_H= \sup \{r : \lim_m \max_{\|\alpha\| = m} |c_\alpha|r^m=0\}$. We note that this definition agrees with the one variable case when $n=1$. If $\rho_H>0$ then $H$ converges to an analytic function on $B_n(\rho_H)$. For each $j\ge 0$ let $P_{H,j}(\xx)=\sum_{\|\alpha\|=j} c_\alpha \xx^\alpha$; this is the homogeneous degree $j$ part of $H$. For $\uu\in K^n$ with $|\uu|=1$ we define
$$
H_\uu(z)=H(z\cdot \uu)=\sum_{j=0}^\infty P_{H,j}(\uu) \cdot z^j\in K[[z]].
$$
It follows that $\rho_{H_\uu}\ge \rho_H$ for each $\uu\in K^n$ with $|\uu|=1$.
%%
%%
\begin{lemma}[Multivariable zero estimate]\label{LemmaMVzeros} Let $H(\xx)=\sum_\alpha c_\alpha \xx^\alpha\in K[[\xx]]
\source{chabauty-coleman-bound-surfaces}
\notready$ and let $\uu\in K^n$ with $|\uu|=1$. Suppose that there is an integer  $N\ge 1$ and a real number $M\ge 1$ such that
\source{chabauty-coleman-bound-surfaces}
\begin{itemize}
\item[(i)] For each $\alpha$ we have $|c_\alpha|\le M^{\|\alpha\|-1}$.
\item[(ii)] $\max\{|P_{H,j}(\uu)| : 0\le j\le N\}\ge 1$ (in particular, $H_\uu\ne 0$). 
\end{itemize}
Then $\rho_H\ge M^{-1}$ and for every $0<r<M^{-1}$ we have
$$
\nn_0(H_\uu,r)\le \left(N-\frac{\log M}{\log (r^{-1})}\right)\left(1-\frac{\log M}{\log (r^{-1})}\right)^{-1}.
$$
\end{lemma}
%%
%%
\begin{proof} 
By (i), $\rho_H\ge M^{-1}$. Fix $0<r<M^{-1}$, in particular $r<1$. From (ii) and $r<1$ we get
$$
r^N\le \max_{0\le j\le N} |P_{H,j}(\uu)|r^j.
$$
Let $\lambda = (\log M)/\log(r^{-1})$ and note that $0<\lambda<1$. For each $m> (N-\lambda)/(1-\lambda)$ we have
$$
M^{m-1}r^m=M^{-1}(Mr)^m<M^{-1}(Mr)^{(N-\lambda)/(1-\lambda)}=r^N
$$
because $Mr<1$. By (i) and these observations, for each $m> (N-\lambda)/(1-\lambda)$ we get
\begin{equation}\label{EqnZeros1}
\source{chabauty-coleman-bound-surfaces}
|P_{H,m}(\uu) |r^m \le M^{m-1} r^m  < r^N \le \max_{0\le j\le N} |P_{H,j}(\uu )|r^j.
\end{equation}
Let $N'=\lfloor (N-\lambda)/(1-\lambda)\rfloor$ and note that $N'\ge N$.  From \eqref{EqnZeros1} we deduce 
$$
|H_{\uu}|_r\le \max\{|P_{H,j}(\uu )|r^j : 0\le j\le N'\}\quad \mbox{ and }\quad \nu(H_\uu, r)\le N'.
$$
The result follows from Lemma \ref{LemmaZeros1var}.
\end{proof}
%%%%%%%%%%%%%%%%%%%%%%%%%%%%%%%%%%%%%%
%%%%%%%%%%%%%%%%%%%%%%%%%%%%%%%%%%%%%%
%%%%%%%%%%%%%%%%%%%%%%%%%%%%%%%%%%%%%%
%%%%%%%%%%%%%%%%%%%%%%%%%%%%%%%%%%%%%%
%%%%%%%%%%%%%%%%%%%%%%%%%%%%%%%%%%%%%%
%%%%%%%%%%%%%%%%%%%%%%%%%%%%%%%%%%%%%%
